\documentclass[10pt]{article}

\usepackage[a4paper, total={6in, 9.5in}]{geometry}
\usepackage[dvipsnames,svgnames]{xcolor}
\usepackage[shortlabels]{enumitem}
\usepackage[framemethod=TikZ]{mdframed}
\usepackage{amsmath,amssymb,amsthm}
\usepackage[colorlinks]{hyperref}

\hypersetup{
	colorlinks=true,
	linkcolor=blue,
	filecolor=blue,      
	urlcolor=blue,
	pdftitle={MAT344 Notes}
}

\usepackage{mathtools}
\usepackage{thmtools}
\usepackage[textsize=scriptsize]{todonotes}

\renewcommand{\tilde}{\widetilde}
\renewcommand{\hat}{\widehat}
\renewcommand{\setminus}{\!-\!}
\providecommand{\ol}{\overline}
\providecommand{\eps}{\varepsilon}
\providecommand{\half}{\frac{1}{2}}
\providecommand{\dang}{\measuredangle} %% Directed angle
\providecommand{\dd}{\mathrm d}
\providecommand{\CC}{\mathbb C}
\providecommand{\FF}{\mathbb F}
\providecommand{\HH}{\mathbb H}
\providecommand{\NN}{\mathbb N}
\providecommand{\QQ}{\mathbb Q}
\providecommand{\RR}{\mathbb R}
\providecommand{\ZZ}{\mathbb Z}
\providecommand{\RP}{\mathbb RP}
\providecommand{\CP}{\mathbb CP}
\providecommand{\dg}{^\circ}
\providecommand{\ind}[1]{\mathbf 1_{#1}}
\providecommand{\ii}{\item}
\providecommand{\alert}[1]{{\sffamily\textbf{\textcolor{magenta}{#1}}}}
\providecommand{\opname}{\operatorname}
\providecommand{\li}{\opname{li}}
\providecommand{\ls}{\opname{ls}}
\providecommand{\supp}{\opname{supp}}
\providecommand{\res}{\opname{Res}}
\providecommand{\inv}{^{-1}}
\providecommand{\setdiff}{\smallsetminus}
\providecommand{\mb}[1]{\mathbf{#1}}
\providecommand{\abs}[1]{\left\lvert{#1}\right\rvert}
\providecommand{\smabs}[1]{\lvert{#1}\rvert}
\providecommand{\norm}[1]{\left\|{#1}\right\|}
\providecommand{\dif}{\;d}
\providecommand{\id}{\operatorname{id}}

\reversemarginpar
\setlength{\marginparwidth}{1in}

\mdfdefinestyle{mdbluebox}{roundcorner=10pt,innerbottommargin=9pt,
	linecolor=blue,backgroundcolor=TealBlue!5,}
\declaretheoremstyle[headfont=\sffamily\bfseries\color{MidnightBlue},
mdframed={style=mdbluebox},]{thmbluebox}

\mdfdefinestyle{mdredbox}{frametitlefont=\bfseries,innerbottommargin=8pt,
	nobreak=true,backgroundcolor=Salmon!5,linecolor=RawSienna,}
\declaretheoremstyle[headfont=\bfseries\color{RawSienna},
mdframed={style=mdredbox},headpunct={\\[3pt]},postheadspace=0pt,]{thmredbox}

\mdfdefinestyle{mdgreenbox}{linecolor=ForestGreen,backgroundcolor=ForestGreen!5,
	linewidth=2pt,rightline=false,leftline=true,topline=false,bottomline=false,}
\declaretheoremstyle[headfont=\bfseries\sffamily\color{ForestGreen!70!black},
mdframed={style=mdgreenbox},headpunct={ --- },]{thmgreenbox}

\mdfdefinestyle{mdorangebox}{linecolor=RedOrange,backgroundcolor=RedOrange!5,
	linewidth=2pt,rightline=false,leftline=true,topline=false,bottomline=false,}
\declaretheoremstyle[headfont=\bfseries\sffamily\color{RedOrange!70!black},
mdframed={style=mdorangebox},headpunct={ --- },]{thmorangebox}

\mdfdefinestyle{mdblackbox}{linecolor=black,backgroundcolor=RedViolet!5!gray!5,
	linewidth=3pt,rightline=false,leftline=true,topline=false,bottomline=false,}
\declaretheoremstyle[mdframed={style=mdblackbox}]{thmblackbox}

\declaretheoremstyle[spaceabove=6pt,spacebelow=6pt]{boldhead}

\declaretheorem[style=thmbluebox,name=Theorem,numberwithin=section]{theorem}

\declaretheorem[style=boldhead,name=Example,sibling=theorem]{example}

\declaretheorem[style=thmbluebox,name=Corollary,sibling=theorem]{corollary}
\declaretheorem[style=thmbluebox,name=Proposition,sibling=theorem]{prop}
\declaretheorem[style=thmbluebox,name=Lemma,sibling=theorem]{lemma}
\declaretheorem[style=thmbluebox,name=Theorem,numbered=no]{theorem*}
\declaretheorem[style=thmbluebox,name=Lemma,numbered=no]{lemma*}
\declaretheorem[style=thmgreenbox,name=Claim,numberwithin=theorem]{claim}
\declaretheorem[style=thmgreenbox,name=Claim,numbered=no]{claim*}
\declaretheorem[style=boldhead,name=Remark,sibling=theorem]{remark}
\declaretheorem[style=boldhead,name=Remark,numbered=no]{remark*}
\declaretheorem[style=thmgreenbox,name=Definition,sibling=theorem]{definition}
\declaretheorem[style=thmgreenbox,name=Definition,numbered=no]{definition*}
\declaretheorem[style=thmorangebox,name=Warning,sibling=theorem]{warning}
\declaretheorem[style=thmorangebox,name=Warning,numbered=no]{warning*}
\declaretheorem[style=thmorangebox,name=Note,numbered=no]{note}
\declaretheorem[style=thmblackbox,name=Exercise,sibling=theorem]{exercise}
\declaretheorem[style=thmblackbox,name=Exercise,numbered=no]{exercise*}

\newenvironment{soln}{\begin{proof}[Solution]}{\end{proof}}
\newenvironment{walkthrough}{\noindent \textbf{Walkthrough.}}{\medskip}
\newlist{walk}{enumerate}{3}
\setlist[walk]{label=\bfseries (\alph*)}

\providecommand{\pow}{\operatorname{Pow}}
%\providecommand{\vocab}[1]{{{\sffamily\textolor{ForestGreen}{#1}}}}
\providecommand{\vocab}[1]{{{\sffamily\textit{#1}}}}
\providecommand{\lcm}{\operatorname{lcm}}
\providecommand{\ord}{\operatorname{ord}}
\providecommand{\Int}{\operatorname{Int}}
\providecommand{\rk}{\operatorname{rk}}
\providecommand{\Mat}{\operatorname{Mat}}

\usepackage{asymptote}
\begin{asydef}
	defaultpen(fontsize(10pt));
	unitsize(1cm);
	size(8cm); // set a reasonable default
	usepackage("amsmath");
	usepackage("amssymb");
	settings.tex="pdflatex";
	settings.outformat="pdf";
	// Replacement for olympiad+cse5 which is not standard
	import geometry;
	// recalibrate fill and filldraw for conics
	void filldraw(picture pic = currentpicture, conic g, pen fillpen=defaultpen, pen drawpen=defaultpen)
	{ filldraw(pic, (path) g, fillpen, drawpen); }
	void fill(picture pic = currentpicture, conic g, pen p=defaultpen)
	{ filldraw(pic, (path) g, p); }
	// some geometry
	pair foot(pair P, pair A, pair B) { return foot(triangle(A,B,P).VC); }
	pair orthocenter(pair A, pair B, pair C) { return orthocentercenter(A,B,C); }
	pair centroid(pair A, pair B, pair C) { return (A+B+C)/3; }
	// cse5 abbreviations
	path CP(pair P, pair A) { return circle(P, abs(A-P)); }
	path CR(pair P, real r) { return circle(P, r); }
	pair IP(path p, path q) { return intersectionpoints(p,q)[0]; }
	pair OP(path p, path q) { return intersectionpoints(p,q)[1]; }
	path Line(pair A, pair B, real a=0.6, real b=a) { return (a*(A-B)+A)--(b*(B-A)+B); }
	// cse5 more useful functions
	picture CC() {
		picture p=rotate(0)*currentpicture;
		currentpicture.erase();
		return p;
	}
	pair MP(Label s, pair A, pair B = plain.S, pen p = defaultpen) {
		Label L = s;
		L.s = "$"+s.s+"$";
		label(L, A, B, p);
		return A;
	}
	pair Drawing(Label s = "", pair A, pair B = plain.S, pen p = defaultpen) {
		dot(MP(s, A, B, p), p);
		return A;
	}
	path Drawing(path g, pen p = defaultpen, arrowbar ar = None) {
		draw(g, p, ar);
		return g;
	}
\end{asydef}


\title{\vspace{-2.0cm}MAT344 Notes}
\author{\large Aditya Pahuja}
\date{\large \today}
\begin{document}
\maketitle
\tableofcontents
\pagebreak

\section{Differentiability in $\CC$}

\setcounter{subsection}{-1}
\subsection{Precalculus review}

\begin{definition}[$\CC$]
    \label{def:complex-numbers}
    The \vocab{complex numbers} $\CC$ is the direct sum
    $\RR\oplus i\RR$, where $i\equiv(0,1)$, $1\equiv(1,0)$;
    it is given the structure of a commutative,
    associative, unital $\RR$-algebra with the multiplication
    \begin{equation*}
	(a,b)\cdot(c,d) = (ac-bd,ad+bc)
    \end{equation*}
    (so in particular $i^2 = -1$).
\end{definition}

\begin{exercise}
    A list of properties you should remember from high school:
    \begin{enumerate}[(a)]
        \ii Show that $\CC$ is a field, in with
	$(a+bi)\inv = \frac{a-bi}{a^2+b^2}$.
	\ii Show that complex conjugation, defined as
	$\ol{a+bi}\equiv a-bi$, is a field automorphism of $\CC$
	that fixes $\RR$.
	\ii Show that the map $\abs{\bullet}\colon\CC\to\RR_+$
	given by $\abs{a+bi} = \sqrt{a^2+b^2}$
	(equivalently, $\abs{z}^2 = z\ol z)$)
	is a norm on $\CC$.
	\ii Show that any $z\in\CC\setminus\{0\}$
	can be written uniquely as
	$re^{i\theta} = r(\cos\theta + i\sin\theta)$
	for $r\in\RR_+$ and $\theta\in\RR/2\pi\ZZ$
	(here, the exponential map is defined via its power series).
	In particular, this suggests that multiplication by $z=re^{i\theta}$
	corresponds to scaling about the origin by $r$
	and rotating counterclockwise by $\theta$.
    \end{enumerate}
\end{exercise}

Property (d) in particular suggests why $\CC$ is so nice:
multiplication by a complex number is much more rigid than
any old linear transformation on $\RR^2$, as we will see shortly.

\subsection{Holomorphic functions}

Recall that a map $f\colon\RR^n\to\RR^n$ is differentiable
if it locally behaves like an $\RR^n$-linear map.
Naturally, one gets the following definition for complex differentiability:
\begin{definition}[Holomorphic]
    \label{def:holomorphic}
    A function $f\colon U\to\CC$ defined on an open subset $U\subset\CC$
    is \vocab{holomorphic} or \vocab{complex analytic} at $z_0\in U$ if
    \begin{equation*}
	f'(z_0)\equiv\lim_{h\to 0}\frac{f(z_0+h) - f(z_0)}h
    \end{equation*}
    exists in $\CC$ (i.e. the limit is taken over $h\in\CC$ tending to $0$).
\end{definition}

As with functions on $\RR$, sums, products, quotients (where defined),
and compositions (where defined) of holomorphic functions
are also holomorphic.

\subsubsection{Conformal linear maps, Cauchy-Riemann equations}
Observe that a $\CC$-linear map $L\colon\CC\to\CC$ is
uniquely determined by $L(1)$, since $L(z) = zL(1)$ by linearity;
thus any $\CC$-linear map corresponds to multiplying by
a fixed complex number $a+bi$. As an $\RR$-linear map,
multiplication by $a+bi$ can be expressed as the matrix
\begin{equation*}
    \begin{pmatrix}
	a&-b\\b&a
    \end{pmatrix}
\end{equation*}
with respect to the standard basis ($1$ and $i$) of $\CC$.

As mentioned at the end of section 1.0, such linear transformations
are angle- and orientation-preserving (equivalently,
they preserve directed angles --- that is, if $\ell_1$ and $\ell_2$
are lines such that the counterclockwise angle
from $\ell_1$ to $\ell_2$ is $\theta$, then the counterclockwise angle
from $L(\ell_1)$ to $L(\ell_2)$ is also $\theta$).
Such a linear map is called \vocab{conformal}.

\begin{remark}
    More generally, we call a differentiable map $f\colon\RR^2\to\RR^2$
    (or between two Riemannian manifolds of dimension $2$)
    conformal if it preserves angles between curves,
    where the angle between two curves at a given point of intersection $p$
    is the angle between the tangent lines at that point;
    this is equivalent to $\dd f_p$ being a conformal linear map.
\end{remark}

\begin{prop}
    \label{prop:conformal-is-complex-mult}
    All conformal $\RR$-linear maps $L\colon\CC\to\CC$ are $\CC$-linear.
\end{prop}

\begin{proof}
    Let $z\in\CC$. The product $iz$ is $z$ rotated
    $90^\circ$ counterclockwise about the origin,
    so $L(iz)$ is the counterclockwise $90^\circ$ rotation
    of $L(z)$ about the origin, i.e. $L(iz) = iL(z)$.
    Thus, by $\RR$-linearity,
    \begin{equation*}
        L((a+bi)z) = aL(z) + bL(iz) = (a+bi)L(z)
    \end{equation*}
    for any $a,b\in\RR$, which means $L$ is $\CC$-linear.
\end{proof}

Let $f\colon\CC\to\CC$ be a holomorphic function.
We can consider $f$ as a map on $\RR^2$ by treating it as
a function of the real and imaginary parts;
in this context, the derivative of $f(x,y) = (f_1(x,y),f_2(x,y))$
(where $f_1,f_2\colon\RR^2\to\RR$) is the linear transformation
\begin{equation*}
    \dd f = \begin{pmatrix}
        \partial f_1/\partial x &\partial f_1/\partial y\\
        \partial f_2/\partial x &\partial f_2/\partial y
    \end{pmatrix}.
\end{equation*}
The limit defining the derivative of $f$
corresponds to this linear transformation,
so $\dd f$ must be a conformal linear map; this means that we must have
\begin{equation*}
    \frac{\partial f_1}{\partial x} = \frac{\partial f_2}{\partial y},
    \qquad \frac{\partial f_1}{\partial y} = -\frac{\partial f_2}{\partial x}
\end{equation*}
at any point at which $f$ is holomorphic.
This system of equations is called the \vocab{Cauchy-Riemann equations}.
In fact, a partial converse holds too: the Cauchy-Riemann equations
exactly determine whether $f$ is complex differentiable at a given point
if $f$ is sufficiently regular.

\begin{prop}
    \label{prop:CR-eqns-iff-holomorphic}
    Let $f\colon U\to\CC$, with $U\subseteq\CC$ an open subset,
    have continuous partial derivatives on all of $U$
    (we say that such a function $f$ is $C^1$,
    since it is $C^1$ as a function on (a subset of) $\RR^2$).
    Then, $f$ is holomorphic if and only if
    the Cauchy-Riemann equations hold for all $z_0\in U$.
\end{prop}

\begin{proof}
    If $f$ is holomorphic, then the limits
    \begin{align*}
	\lim_{\substack{h\to 0\\h\in\RR}} \frac{f(z_0+h) - f(z_0)}h
	&= \frac{\partial f}{\partial x}(z_0)
	= \frac{\partial f_1}{\partial x}(z_0)
	+ i\frac{\partial f_2}{\partial x}(z_0)\\
	\lim_{\substack{h\to 0\\h\in i\RR}} \frac{f(z_0+h) - f(z_0)}{h}
	&= \frac 1i\frac{\partial f}{\partial y}(z_0)
	= -i\frac{\partial f_1}{\partial y}(z_0)
	+ \frac{\partial f_2}{\partial y}(z_0)
    \end{align*}
    are equal. Comparing the real and imaginary parts
    yields the Cauchy-Riemann equations.

    Conversely, if $f$ has continuous partial derivatives,
    then it is differentiable as a map on (a subset of) $\RR^2$.
    This derivative obeys the Cauchy-Riemann equations,
    so it is a $\CC$-linear map, say corresponding to multiplication by
    $z\in\CC$. Thus, by definition of differentiability in $\RR^n$,
    \begin{equation*}
	0 = \lim_{h\to 0}\frac{f(z_0+h) - f(z_0) - \dd f_{z_0}h}h
	= \lim_{h\to 0}\frac{f(z_0+h) - f(z_0) - zh}h
    \end{equation*}
    so the limit defining $f'(z_0)$ exists.
\end{proof}

\begin{remark}
    [Looman-Menchoff theorem]
    In fact, any continuous function $f\colon U\to \CC$
    whose partials exist and satisfy the Cauchy-Riemann equations
    on all of $U$ is holomorphic on $U$;
    this is called the \vocab{Looman-Menchoff theorem}.
    The continuity assumption is required because
    $f(z) = \exp(-z^{-4})$ has partials that exist and satisfy
    Cauchy-Riemann everywhere, but is not continuous
    (hence not holomorphic) at zero.
    Additionally, if $f$ is only required to have partials
    satisfying the Cauchy-Riemann equations at a single point,
    then $f\colon\CC\to\CC$ defined as $f(0) = 0$ and
    $f(z) = z^5/\abs{z}^4$ for $z\ne 0$ is continuous at $0$
    but not holomorphic at $0$, since
    \begin{equation*}
	\lim_{z\to 0} \frac{f(z) - 0}{z}
	= \lim_{z\to 0} \frac{z^4}{\abs z^4}
    \end{equation*}
    is easily seen to not exist, as taking the limit along the real line
    and taking the limit along the line $\opname{Re}z = \opname{Im}z$
    respectively give $1$ and $-1$.
\end{remark}

\subsubsection{The operators $\frac{\partial}{\partial z}$
and $\frac{\partial}{\partial\ol z}$}
\begin{definition}[$\frac{\partial}{\partial z}$
    and $\frac{\partial}{\partial\ol z}$]
    \label{def:dz-dzbar}
    The formal operators $\frac{\partial}{\partial z}$
    and $\frac{\partial}{\partial\ol z}$ are defined as
    \begin{equation*}
	\frac{\partial}{\partial z}
	\equiv \half\left(\frac{\partial}{\partial x}
	- i\frac{\partial}{\partial y}\right),\qquad
	\frac{\partial}{\partial \ol z}
	\equiv \half\left(\frac{\partial}{\partial x}
	+ i\frac{\partial}{\partial y}\right).
    \end{equation*}
\end{definition}

\begin{prop}
    \label{prop:dzbar-zero-iff-holo}
    A $C^1$ function $f\colon U\to\CC$
    (where, as usual, $U\subseteq\CC$ is open)
    is holomorphic if and only if $\frac\partial{\partial\ol z}f = 0$
    on all of $U$.
\end{prop}

\begin{proof}
    Write $f = f_1 + if_2$ for real-valued functions $f_1$ and $f_2$.
    We can manually compute
    \begin{equation*}
	2\frac\partial{\partial\ol z}f
	= \frac{\partial f_1}{\partial x} - \frac{\partial f_2}{\partial y}
	+ i\left(\frac{\partial f_1}{\partial y}
	+ \frac{\partial f_2}{\partial x}\right),
    \end{equation*}
    which is zero if and only if the real and imaginary parts are zero,
    which is equivalent to the Cauchy-Riemann equations.
    The statement then follows from
    Proposition \ref{prop:CR-eqns-iff-holomorphic}.
\end{proof}

\begin{remark}
    We can compute
    \begin{equation*}
	\frac\partial{\partial z}\frac\partial{\partial\ol z}
	= \frac\partial{\partial\ol z}\frac\partial{\partial z}
	= \frac 14\left(\frac{\partial^2}{\partial^2 x}
	+ \frac{\partial^2}{\partial^2y}\right) = \frac 14\Delta
    \end{equation*}
    where $\Delta = \frac{\partial^2}{\partial^2 x}
    + \frac{\partial^2}{\partial^2y}$ is called the \vocab{Laplacian} of $f$.
\end{remark}

\begin{prop}
    \label{prop:dz-equals-derivative}
    Let $f\colon U\to\CC$ be holomorphic at $z_0\in U$. Then
    \begin{equation*}
	\frac\partial{\partial z}f(z_0) = f'(z_0).
    \end{equation*}
\end{prop}

\begin{proof}
    Write $f = f_1 + if_2$ for real-valued functions $f_1,f_2\colon U\to\RR$.
    Then,
    \begin{equation*}
	\frac\partial{\partial z}f
	= \half\left(\frac{\partial f_1}{\partial x}
	+ \frac{\partial f_2}{\partial y}\right)
	+ \frac i2\left(\frac{\partial f_2}{\partial x}
	- \frac{\partial f_1}{\partial y}\right).
    \end{equation*}
    Since $f$ is holomorphic at $z_0$,
    the Cauchy-Riemann equations imply that
    multiplication by $\frac\partial{\partial z}f(z_0)$
    corresponds to the linear transformation
    \begin{equation*}
        \begin{pmatrix}
	    \partial f_1/\partial x & \partial f_1/\partial y\\
	    \partial f_2/\partial x & \partial f_2/\partial y
        \end{pmatrix}
    \end{equation*}
    \todo{i confuse myself}
\end{proof}

\subsection{Functions on Riemann surfaces}

Let $\CC_{\infty} = \CC\cup\{\infty\}$ be the one-point compactification
of $\CC$ (which is homeomorphic to $S^2$).
Let $U = \CC_{\infty} - \{\infty\} = \CC$ and $V = \CC_\infty-\{0\}$,
and define
\begin{align*}
    \varphi\colon U\to\CC,\quad f(z) &= z,\\
    \psi\colon V\to\CC,\quad f(z) &= \frac 1z
\end{align*}
with the convention $\frac 1\infty = 0$.
These are both homeomorphisms (since they continuous
and are both their own inverse) and the transition map
\begin{equation*}
    \psi\circ\varphi\inv\colon\varphi(U\cap V)\to\psi(U\cap V),
    \quad \psi\circ\varphi\inv(z) = \frac 1z
\end{equation*}
is a \vocab{biholomorphism} (holomorphic map with holomorphic inverse),
so the charts $(U,\varphi)$ and $(V,\psi)$ turn $\CC_\infty$
into a \vocab{complex $1$-manifold}, also called a \vocab{Riemann surface}.
The Riemann surface $\CC_\infty$ is known as the \vocab{Riemann sphere}.

Concretely, the idea of the maps $\varphi$ and $\psi$ is that
they describe how $\CC_\infty$ looks like some copies of $\CC$ glued together,
with $\varphi\circ\varphi\inv$ being a biholomorphism meaning
that the way they are glued is compatible with the kind of
structure on $\CC$ that we care about.
Since being holomorphic only depends on the behavior of a function locally,
we can use these charts to define holomorphic functions
between any two Riemann surfaces.

\todo{define riemann surface properly}

\todo{define holomorphic map between riemann surfaces}

\begin{example}
    [Automorphisms of the Riemann sphere]
    Let $a,b,c\in\CC$ with $a\ne b$ and $c\ne 0$.
    The function $f\colon\CC_\infty\to\CC_\infty$
    defined by $f(z) = c\cdot\frac{z-a}{z-b}$
    when $z\notin\{b,\infty\}$, $f(b) = \infty$, and $f(\infty) = c$
    is a biholomorphism of $\CC_\infty$;
    in fact, these are the only biholomorphisms on $\CC_\infty$.
    \todo{idk difficulty of pf}
\end{example}

\begin{example}
    Let $p,q\colon\CC\to\CC_\infty$ be complex polynomials
    with no common factors. Define
    \begin{equation*}
        f(z) = \begin{cases}
	    p(z)/q(z) & q(z)\ne 0\\
	    \infty & q(z) = 0.
        \end{cases}
    \end{equation*}
    Since $f$ is holomorphic whenever $q(z)\ne 0$
    and $\frac 1f$ is holomorphic whenever $p(z)\ne 0$,
    this makes $f\colon\CC\to\CC_\infty$ a holomorphic function.
    If $a\in\CC$ is a point at which $f(z) = \infty$,
    there exists a positive integer $n$ such that
    \begin{equation*}
	f(z) = \frac{1}{(z-a)^n}g(z)
    \end{equation*}
    for holomorphic $g$ that is nonzero on (a neighborhood of) $a$.
    The complex number $a$ is called a \vocab{pole} of $f$,
    and the integer $n$ is called its \vocab{order}.
    Similarly, if $f(b) = 0$, then
    $f(z) = (z-b)^nh(z)$ for $h$ that is nonzero and holomorphic
    in a neighborhood of $b$; in this case $b$ is a \vocab{zero} of $f$
    with \vocab{order} $n$.
    Note that a pole of $f$ with order $n$ is a zero of $\frac 1f$
    with order $n$, and vice versa.
\end{example}

\todo{give another example of riemann surface using log}


\todo{write down power series facts}
\section{Differential geometry crash course}
\begin{definition}[Tangent vector, cotangent vector]
    \label{def:tangent-cotangent}
    todo
\end{definition}

\todo{induced maps on tangent/cotangent spaces via a smooth map $\RR^m\to\RR^n$}

\begin{definition}[$1$-form]
    \label{def:1-form}
    todo
\end{definition}

\begin{definition}[Integral of $1$-form]
    \label{def:line-integral}
    todo
\end{definition}

\todo{good domains, orientation, Green's theorem}

\pagebreak

\section{The Cauchy integral formula and consequences}
\todo{write down important results from HWs}

\todo{restructure sections}
Now that we have established the basic geometric and algebraic terminology
and some facts about $\CC$ we can actually do some complex analysis.

\subsection{The Cauchy integral formula}

\todo{morera}
\begin{theorem}[Cauchy's integral theorem]
    \label{thm:cauchy-integral-theorem}
    Let $\omega = f(z)\dd z$ be a holomorphic $1$-form
    on a good domain $\Omega$. Then
    \begin{equation*}
	\int_{\partial\Omega}\omega = 0.
    \end{equation*}
\end{theorem}

\begin{proof}
    Apply Green's theorem to get
    \begin{equation*}
	\int_{\partial\Omega}\omega = \int_{\Omega}\dd\omega = 0
    \end{equation*}
    where the integral is zero because $\dd\omega = 0$, as $f$ is holomorphic.
\end{proof}

\begin{theorem}[Cauchy's integral formula]
    \label{thm:cauchy-integral-formula}
    Let $\Omega$ be a good domain and $f\in\mathcal O(\Omega)$.
    For every $z\in\Omega$,
    \begin{equation*}
	f(z) =
	\frac 1{2\pi i}\int_{\partial\Omega}\frac{f(\zeta)}{\zeta-z}\dd\zeta.
    \end{equation*}
\end{theorem}

\begin{proof}
    Fix $z\in\Omega$ and let $\eps>0$ be such that
    the closed ball $\ol B_\eps(z)$ is contained in $\Omega$.
    Since $\frac{f(\zeta)}{\zeta - z}$ is a holomorphic function of $\zeta$
    on $\Omega - B_\eps(z)$, Green's theorem shows that
    \begin{equation*}
	\int_{\partial\Omega}\frac{f(\zeta)}{\zeta-z}\dd\zeta
	= \int_{\abs{\zeta-z} = \eps} \frac{f(\zeta)}{\zeta-z}\dd\zeta.
    \end{equation*}
    Then, parametrizing the circle in the obvious way, we get
    \begin{equation*}
	\int_{\abs{\zeta-z} = \eps} \frac{f(\zeta)}{\zeta-z}\dd\zeta.
	= \int_0^{2\pi}\frac{f(z+\eps e^{i\theta})}{\eps e^{i\theta}}
	\cdot \eps ie^{i\theta}\dd\theta
	= i\int_0^{2\pi}f(z+\eps e^{i\theta})\dd\theta.
    \end{equation*}
    The equation we got from Green's theorem tells us that
    this integral is independent of $\eps$,
    so it is equal to its limit as $\eps$ goes to zero;
    moreover, we can compute that
    \begin{equation*}
	\abs{\int_0^{2\pi}f(z+\eps e^{i\theta}) - f(z)\dd\theta}
	< \int_0^{2\pi} C\eps\dd\theta = 2\pi C\eps,
    \end{equation*}
    where $C$ is a constant such that
    $\abs{f(z + \eps e^{i\theta}) - f(z)} < C\abs{\eps e^{i\theta}}$
    for sufficiently small $\eps$, which exists because $f$ is differentiable
    (namely take any $C > \abs{f'(z)}$). In summary,
    \begin{align*}
	\int_{\partial\Omega}\frac{f(\zeta)}{\zeta - z}\dd\zeta
	&= \int_{\abs{\zeta - z} = \eps}\frac{f(\zeta)}{\zeta - z}\dd\zeta\\
	&= i\int_0^{2\pi}f(z + \eps e^{i\theta})\dd\theta
	= i\cdot \lim_{\eps\to 0}
	\int_0^{2\pi}f(z + \eps e^{i\theta})\dd\theta\\
	&= i\int_0^{2\pi} f(z)\dd\theta = 2\pi if(z).\qedhere
    \end{align*}
\end{proof}

\pagebreak
\begin{theorem}[Holomorphic functions are analytic]
    \label{thm:holomorphic-implies-analytic}
    Let $\Omega$ be a good domain and $f\in\mathcal O(\Omega)$.
    For every $z_0\in\Omega$, $f$ has a convergent power series
    in a ball around $z_0$.
\end{theorem}

\begin{proof}
    Let $r$ be the distance from $z_0$ to $\partial\Omega$
    and suppose that $z\in\ol B_{r_0}(z_0)$ for some $r_0 < r$.
    Then
    \begin{align*}
	f(z)
	&= \frac{1}{2\pi i}\int_{\partial\Omega}
	\frac{f(\zeta)}{\zeta-z}\dd\zeta
	= \frac 1{2\pi i}\int_{\partial\Omega}
	\frac{f(\zeta)}{(\zeta - z_0) - (z-z_0)}\\
	&= \frac 1{2\pi i}\int_{\partial\Omega}
	\left(\frac{f(\zeta)}{\zeta - z_0}\sum_{n=0}^\infty
	\left(\frac{z-z_0}{\zeta-z_0}\right)^n\right)\dd\zeta\\
	&= \sum_{n=0}^\infty \left( \frac 1{2\pi i}
	\int_{\partial\Omega}\frac{f(\zeta)}{(\zeta-z_0)^{n+1}}\dd\zeta
        \cdot (z-z_0)^n\right)
    \end{align*}
    where we can expand the geometric series beacuse
    $\abs{z-z_0} < \abs{\zeta - z_0}$ by assumption
    and we can pull the sum out of the integral because
    the geometric series converges uniformly
    over $z\in\ol B_{r_0}(z_0)$ (say, by the Weierstrass $M$-test).
    This is thus a power series expansion for $f(z)$
    over all $z\in B_r(z_0)$.
\end{proof}

\begin{prop}
    [Analytic continuation]
    \label{prop:analytic-continuation}
    Let $U\subseteq\CC$ be open and connected
    and suppose that $f\in\mathcal O(U)$.
    \begin{enumerate}[(1)]
        \ii $f$ is infinitely differentiable.
	\ii The zeroes of $f$ are isolated unless $f$ is
	the zero function on $U$.
	\ii If $f$ and $g$ are holomorphic functions on $U$
	that agree on any open set, then they are the same function.
    \end{enumerate}
\end{prop}

\begin{proof}
    (1) This is immediate from Theorem \ref{thm:holomorphic-implies-analytic}.
    \\

    (2) Suppose $z_0\in U$ satisfies $f(z_0) = 0$.
    Then, in a neighborhood of $z_0$,
    \begin{equation*}
	f(z) = \sum_{n=0}^\infty a_n(z-z_0)^n.
    \end{equation*}
    If there is a $k>0$ such that $a_k \ne 0$,
    take the smallest such $k$, so that
    $f(z) = (z-z_0)^k g(z)$ for $g$ defined in a neighborhood of $z_0$ by
    \begin{equation*}
	g(z) = \sum_{n=0}^\infty a_{n+k}(z-z_0)^n.
    \end{equation*}
    Since $g(z_0) = \lim_{z\to z_0}g(z) \ne 0$,
    $g$, hence $f$, is nonzero in a neighborhood of $z_0$,
    so $z_0$ is an isolated zero of $f$.

    On the other hand, if all the $a_n$ are zero,
    then $f$ is zero in a neighborhood of $z_0$.
    This implies that the set of non-isolated zeroes of $f$
    is open in $U$. At the same time, the set of non-isolated zeroes
    must be closed in $U$, because $f\inv(0)$ is closed
    and a closed set minus some isolated points is still closed.
    Connectedness of $U$ then tells us that
    the set of non-isolated zeroes of $f$ is either
    all of $U$ (in which case $f$ is the zero function)
    or empty (in which case all the zeroes of $f$ are isolated).
    \\

    (3) If $f$ and $g$ agree on an open set, then
    $f-g$ has a non-isolated zero, so $f-g\equiv 0$ by (2).
\end{proof}

\pagebreak
\subsection{Derivatives as integrals}
The fact that holomorphic functions are analytic
and that their coefficients can be computed
via certain integrals gives us the following formula:

\begin{corollary}
    \label{coro:nth-derivative-integral}
    If $f$ is holomorphic on a good domain $\Omega$, then
    \begin{equation*}
	f^{(n)}(z) = \frac{n!}{2\pi i}\int_{\partial\Omega}
	\frac{f(\zeta)}{(\zeta-z)^{n+1}}\dd\zeta.
    \end{equation*}
\end{corollary}

\begin{proof}
    Immediate from differentiating the power series expansion around $z$
    (recall that the derivative of a power series
    is obtained by termwise differentiation).
\end{proof}

% One can also prove Corollary \ref{coro:nth-derivative-integral}
% with the following fact, taking $K = \partial\Omega$
% and $\mu = f(\zeta)\dd s$.\todo{idk what this measure is bro}
% \begin{lemma}
%     \label{prop:borel-measure-lemma}
%     Let $K\subseteq\CC$ be compact and $\mu$ a finite Borel measure on $K$.
%     The function
%     \begin{equation*}
% 	F(z) = \int_K \frac 1{\zeta-z}\dd\mu(\zeta)
%     \end{equation*}
%     is holomorphic on $\CC - K$, with
%     \begin{equation*}
% 	F^{(n)}(z) = n!\int_K \frac 1{(\zeta-z)^{n+1}}\dd\mu(\zeta).
%     \end{equation*}
% \end{lemma}
% 
% \begin{proof}
%     \todo{pf}
% \end{proof}

One reason that this formula is useful because it
lets us easily estimate derivatives; for example:

\begin{lemma}
    \label{lem:derivative-of-cauchy-seq}
    If $\{f_k\}_{k=1}^\infty$ is a Cauchy sequence
    in $\mathcal O(\ol B_r(z_0))$, then, for all $n\ge 1$,
    the sequence of $n$th derivatives $\{f_k^{(n)}\}_{k=1}^\infty$
    is also Cauchy when restricted to $\ol B_{r_0}(z_0)$
    for any $r_0 < r$ (where the space of holomorphic functions
    is endowed with the sup metric).
\end{lemma}

\begin{proof}
    Let $k,\ell\in\NN$ be such that
    $\abs{f_k-f_{\ell}} < \eps$ on $\partial B_r(z_0)$.
    Then by Corollary \ref{coro:nth-derivative-integral},
    \begin{equation*}
	f_k^{(n)}(z) - f_\ell^{(n)}(z)
	= \frac{n!}{2\pi i}\int_{\abs{\zeta-z_0}=r}
	\frac{f_k(\zeta)-f_\ell(\zeta)} {(\zeta-z)^{n+1}}\dd\zeta.
    \end{equation*}
    for all $z\in B_r(z_0)$. Suppose now that $z\in B_{r_0}(z_0)$,
    so that $\abs{\zeta-z}\ge r-r_0$ for $\zeta$ on $\partial B_r(z_0)$.
    We then get
    \begin{align*}
	\smabs{f_k^{(n)}(z) - f_\ell^{(n)}(z)}
	&= \frac{n!}{2\pi}\abs{\int_{\abs{\zeta-z_0}=r}
	\frac{f_k(\zeta)-f_\ell(\zeta)}{(\zeta-z)^{n+1}}\dd z}\\
	&= \frac{n!}{2\pi}\abs{\int_0^{2\pi} \frac{f_k(\zeta)-f_\ell(\zeta)}
	{(\zeta-z)^{n+1}}\cdot ire^{i\theta}\dd \theta}\\
	&\le \frac{n!\cdot r}{2\pi}\int_0^{2\pi}\abs{\frac{f_k(\zeta)-f_\ell(\zeta)}
	{(\zeta-z)^{n+1}}}\dd\theta\\
	&< \frac{n!\cdot r}{2\pi}
	\int_0^{2\pi}\frac{\eps}{(r-r_0)^{n+1}}\dd\theta
	= \frac{n!\cdot r\eps}{(r-r_0)^{n+1}}.
    \end{align*}
    Since the coefficient of $\eps$ depends only on $n$, $r$, and $r_0$,
    which are fixed, this implies that $f_k^{(n)}$ is Cauchy
    in $\ol B_{r_0}(z_0)$ if $f_k$ is Cauchy on $\ol B_r(z_0)$.
\end{proof}

The following is sort of a version of the maximum principle
but for derivatives.
\begin{lemma}
    \label{lem:norm-estimate-nth-derivative}
    Suppose $f$ is holomorphic on $\ol B_r(z_0)$ and $0 < r_0 < r$.
    For all positive integers $n$,
    \begin{equation*}
	\sup_{B_{r_0}(z_0)}\smabs{f^{(n)}}
	\le \frac{n!}{2\pi(r-r_0)^{n+1}}
	\int_{\partial B_r(z_0)}\abs{f(z)}\dd z
	\le \frac{n!\cdot r}{(r-r_0)^{n+1}}\sup_{\partial B_r(z_0)}\abs f
    \end{equation*}
\end{lemma}

\begin{proof}
    By the inequality in the proof of
    Lemma \ref{lem:derivative-of-cauchy-seq}
    with $f$ in place of $f_k$ and the zero function in place of $f_\ell$,
    \begin{equation*}
	\smabs{f^{(n)}(z)}
	\le \frac{n!}{2\pi(r-r_0)^{n+1}}\int_0^{2\pi}
	\abs{f(z_0+re^{i\theta})}r\dd\theta
	\le \frac{n!\cdot r}{(r-r_0)^{n+1}}\sup_{\partial B_r(z_0)}\abs f
    \end{equation*}
    holds for all $z\in B_{r_0}(z_0)$,
    where $\eps$ is the supremum of $\abs f$
    and we have simply not yet estimated $\abs f$ above by $\eps$
    in the middle term. Then, the middle term is easily seen to be
    equal to the integral of $\abs f$ over the boundary of $B_r(z_0)$,
    so taking the supremum over $z\in B_{r_0}(z_0)$
    gives the inequality as stated in the lemma.
\end{proof}

\begin{theorem}[Liouville's theorem]
    \label{thm:liouville}
    If $f$ is \vocab{entire}, meaning it is holomorphic on all of $\CC$,
    and $\abs f$ is bounded, then $f$ is constant.
\end{theorem}

\begin{proof}
    Let $M = \sup_\CC \abs f < \infty$.
    Lemma \ref{lem:norm-estimate-nth-derivative} says
    \begin{equation*}
	\abs{f'(z)} \le \frac{R}{(R-\abs z)^2}\cdot M
    \end{equation*}
    for all $R > \abs{z}$, so taking $R$ arbitrarily large
    implies that $f'(z) = 0$; thus $f'(z) = 0$ everywhere,
    meaning $f$ is constant.
\end{proof}

\begin{theorem}[Fundamental theorem of algebra]
    \label{thm:fundamental-theorem-of-algebra}
    Any nonconstant polynomial in $\CC$ has a root in $\CC$.
\end{theorem}

\begin{proof}
    If $p(z) = a_0 + a_1z + \cdots + a_nz^n \in\CC[z]$
    has no roots, then $\frac 1{p(z)}$ is an entire function. Since
    \begin{equation*}
	\abs{\frac 1{p(z)}} =
	\frac 1{\abs z^n}\cdot\frac 1{\abs{a_n + a_{n-1}z^{-1}
	+ \cdots + a_0z^{-n}}} < \frac 2{\abs{a_nz}^n}
    \end{equation*}
    holds if $\abs z$ is large enough, $\frac 1{p(z)}$
    is bounded outside some ball $\ol B_R(z_0)$.
    On the other hand $\frac 1{p(z)}$ is also bounded inside $\ol B_R(z_0)$
    because $\ol B_R(z_0)$ is compact, so it is bounded on all of $\CC$.
    Liouville's theorem then implies that $\frac 1{p(z)}$,
    hence $p(z)$, must be constant if it has no roots.
\end{proof}

\subsection{Root counting}
The Cauchy integral formula lets us
\emph{count the number of zeroes} of a holomorphic function!
\begin{prop}[Argument principle]
    \label{prop:argument-principle}
    Let $\Omega$ be a good domain and $f\in\mathcal O(\ol\Omega)$.
    If $f|_{\partial\Omega}$ is nowhere zero, then $f$ has
    finitely many, say $N$, zeroes in the interior of $\Omega$, and
    \begin{equation*}
        N = \frac 1{2\pi i}\int_{\partial\Omega}\frac{f'(z)}{f(z)}\dd z.
    \end{equation*}
\end{prop}

\begin{proof}
    The number of zeroes is finite because
    the set of zeroes of $f$ is compact (because it is
    a closed subset of the compact set $\ol\Omega$) and discrete
    (since, by assumption, $f$ is not everywhere zero,
    hence is not the zero function by
    Proposition \ref{prop:analytic-continuation}),
    and compact, discrete sets must be finite.

    If $z_1$, \dots, $z_N$ are the zeroes of $f$
    (where each zero is repeated in this list
    as many times as its multiplicity, e.g. $z_1=z_2 = 0$ if $f(z) = z^2$),
    then $f(z) = (z-z_1)\cdots(z-z_N)g(z)$
    where $g$ is a nonzero holomorphic function on $\ol\Omega$;
    one can verify that $g$ is indeed nonzero on $\ol\Omega$
    by using analyticity. Then it is easy to compute
    \begin{equation*}
	\frac{f'(z)}{f(z)} = \frac{1}{z-z_1} + \cdots + \frac 1{z-z_N}
	+ \frac{g'(z)}{g(z)}
    \end{equation*}
    using the product rule (the trick for computing this quickly is
    to formally differentiate the ``function''
    $\log f = \log(z-z_1) + \cdots + \log(z-z_n) + \log g$).
    When integrating over $\partial\Omega$, the $\frac 1{z-z_k}$ addends
    each become $2\pi i$ by the Cauchy integral formula,
    and since $\frac{g'(z)}{g(z)}$ is holomorphic on $\ol\Omega$,
    its integral is zero by the Cauchy integral theorem, so
    \begin{equation*}
	\frac 1{2\pi i}\int_{\partial\Omega}\frac{f'(z)}{f(z)}\dd z = N.
    \end{equation*}

    Equivalently, one could apply the residue theorem
    (Theorem \ref{thm:residue-theorem}) to $\frac{f'(z)}{f(z)}$.
\end{proof}

Geometrically, this says that the number of solutions to $f(z) = 0$
inside $\Omega$ is the winding number of $f(\partial\Omega)$ around $0$
(as long as $f(\partial\Omega)$ does not pass through zero).

\begin{corollary}
    \label{coro:locally-polynomial}
    Let $U\subseteq\CC$ be open and connected
    and $f\in\mathcal O(U)$ nonconstant.
    For a given $z_0\in U$, suppose that $f(z) - f(z_0) = (z-z_0)^ng(z)$
    where $g\in\mathcal O(U)$ is nonzero at $z_0$
    ($n$ is called the \vocab{order} of the zero $z_0$ of $f$).
    Then, there exists $\delta>0$ such that for all $a\in\CC$
    in $B_\delta(f(z_0))$, the equation $f(z) = a$
    has $n$ solutions in a neighborhood of $z_0$.
\end{corollary}

\begin{proof}
    Let $B_r(z_0)$ be a neighborhood of $z_0$ in which
    $g$ is nonzero, so that
    \begin{equation*}
	\frac 1{2\pi i}\int_{\partial B_r(z_0)} \frac{f'(z)}{f(z)-f(z_0)}\dd z
	= n
    \end{equation*}
    by the argument principle. Then, for $a\in\CC - f(\partial B_r(z_0))$,
    define
    \begin{equation*}
        N(a) \coloneq 
	\frac 1{2\pi i}\int_{\partial B_r(z_0)} \frac{f'(z)}{f(z)- a}\dd z
    \end{equation*}
    to be the number of solutions to $f(z) = a$ inside $B_r(z_0)$.
    Then,
    \begin{align*}
	\abs{N(z_0) - N(a)}
	&= \frac 1{2\pi}\abs{\int_{\partial B_r(z_0)}
	f'(z)\cdot\frac{a-f(z_0)}{(f(z)-a)(f(z)-f(z_0))}\dd z}\\
	&= \frac 1{2\pi}\abs{\int_{0}^{2\pi}f'(z)\cdot
	\frac{a-f(z_0)}{(f(z)-a)(f(z)-f(z_0))}\cdot rie^{i\theta}\dd\theta}\\
	&\le \frac{r\abs{a-f(z_0)}}{2\pi}\int_0^{2\pi}\abs{
	\frac{f'(z)}{(f(z)-a)(f(z)-f(z_0))}}\dd\theta
    \end{align*}
    Since $\partial B_r(z_0)$ is compact and $z_0\notin f(\partial B_r(z_0))$,
    $\abs{f(z)-f(z_0)}$ has a minimum value $M > 0$ on the circle.
    Supposing $\abs{a-z_0} < \delta$ with $\delta < M$,
    we then also have $\abs{f(z-a)} \ge M-\delta$ on the circle. Thus
    \begin{equation*}
	\frac{r\abs{a-f(z_0)}}{2\pi}\int_0^{2\pi}\abs{
	\frac{f'(z)}{(f(z)-a)(f(z)-f(z_0))}}\dd\theta
	\le \frac{r\delta}{M(M-\delta)}\int_0^{2\pi}\abs{f'(z)}\dd\theta,
    \end{equation*}
    and the right-hand side can be made arbitrarily small
    by shrinking $\delta$.
    Thus $N$ is an integer-valued, continuous function
    on its domain, hence it is constant on the component
    of $\CC - f(\partial B_r(z_0))$ containing $z_0$.
    This component is open in $\CC$ because
    it is an open subset of the open set $\CC - f(\partial B_r(z_0))$
    (noting that the image of the circle is compact hence closed),
    so one just needs to pick $\delta$ small enough
    that $B_\delta(f(z_0))$ is contained in this connected component.
\end{proof}

\begin{exercise}[Weak Rouch\'e's theorem]
    \label{exer:rouche-zeroes}
    Prove the following version of Rouch\'e's theorem
    (Theorem \ref{thm:rouches-theorem})
    by appealing to the winding number interpretation
    of the argument principle:
    if $f$ and $g$ are holomorphic functions on good domain $\Omega$
    such that $\abs{f-g} < \abs g$ on $\partial\Omega$,
    then $f$ and $g$ have the same number of zeroes inside $\Omega$.
    The stronger (usual) version of this theorem,
    which we will see later, incorporates some information
    about the poles of $f$ too.
    (Hint: you should be able to construct
    an explicit homotopy $H\colon(\CC-\{0\})\times I\to\CC-\{0\}$
    between $f(\partial\Omega)$ and $g(\partial\Omega)$.)
\end{exercise}

\begin{exercise}
    \label{exer:locally-poly-rouche}
    Use Rouch\'e's theorem to prove Corollary \ref{coro:locally-polynomial}.
\end{exercise}

Two super important corollaries:
\begin{theorem}[Open mapping theorem]
    \label{thm:open-mapping}
    Let $U\subseteq\CC$ be open and connected,
    and $f\in\mathcal O(U)$ nonconstant.
    Then $f$ is an open map.
\end{theorem}

\begin{proof}
    This follows immediately from Corollary
    \ref{coro:locally-polynomial}, since the corollary implies that
    if $a$ is contained in $f(U)$ then so is some open ball around $a$.
\end{proof}

\pagebreak
\begin{theorem}[Maximum modulus principle]
    \label{thm:max-mod-principle}
    Let $\Omega$ be a good domain
    and $f\in\mathcal O(\ol\Omega)$ a holomorphic function.
    For every $z_0$ in the interior of $\Omega$,
    \begin{equation*}
	\abs{f(z_0)} < \sup_{z\in\partial\Omega}\abs{f(z)}
    \end{equation*}
    (note the strict inequality).
\end{theorem}

\begin{proof}
    The open mapping theorem tells us that
    $f$ sends the interior of $\ol\Omega$
    to the interior of $f(\ol\Omega)$.
    Since $\ol\Omega$ is compact,
    $f(\ol\Omega)$ is too, so there is a point in $f(\ol\Omega)$
    maximizing the magnitude over all points in $f(\ol\Omega)$;
    any such point necessarily cannot be in the interior of $f(\ol\Omega)$,
    so neihter can its preimage under $f$.
    Thus any point maximizing the magnitude of $f$
    over $\ol\Omega$ cannot be in the interior of $\Omega$.
\end{proof}

\subsection{Poles}
Let $U\subseteq\CC$ be open and $z_0\in U$.
If $f\in\mathcal O(U - \{z_0\})$, we say that $z_0$ is
an \vocab{isolated singularity} of $f$.
There are three possibilities for the behavior of $f$ around $z_0$.
\begin{enumerate}[(1)]
    \ii If $f$ extends to a holomorphic function on $U$,
    then we call $z_0$ a \vocab{removable singularity} of $f$.
    Such points are characterized by
    $\lim_{z\to z_0}(z-z_0)f(z) = 0$, in which case
    we can define $\tilde f\in\mathcal O(U)$ by taking
    $\tilde f\equiv f$ on $U - \{z_0\}$ and
    \begin{equation*}
	\tilde f(z_0) = \frac 1{2\pi i}\int_{\partial B_\eps(z_0)}
	\frac{f(\zeta)}{\zeta-z_0}\dd\zeta
    \end{equation*}
    (which is of course preordained by the Cauchy integral formula).
\end{enumerate}

\begin{exercise}
    If $f$ is a removable singularity
    (in the sense of the limit definition)
    show that the function $g\colon U\to\CC$ defined by
    $g(z_0) = 0$ and $g(z) = (z-z_0)^2f(z)$ for $z\in U-\{z_0\}$
    is holomorphic on $U$, and that $z_0$
    is a zero of $g$ with order at least $2$.
    Conclude that $f$ admits an extension to $U$.
\end{exercise}

\begin{enumerate}[(1),resume]
    \ii It could also happen that there is a positive integer $n$
    such that $g(z) = (z-z_0)^nf(z)$ extends to a
    nonzero holomorphic function in a neighborhood of $z_0$.
    In this case, $z_0$ is a \vocab{pole} of $f$
    with \vocab{order} $n$. Recall that, dually, $z_0$ is a zero
    of $f$ with order $n$ if $f(z) = (z-z_0)^nh(z)$ for some
    holomorphic $h$ that is nonzero in a neighborhood of $z_0$ ---
    poles and zeroes thus behave ``opposite'' to each other
    (cf. Proposition \ref{prop:argument-principle-poles}
    and Theorem \ref{thm:rouches-theorem}).
    A holomorphic function whose isolated singularities are all poles
    (or removable, but we can just remove those)
    is called \vocab{meromorphic};
    these functions can be thought of as holomorphic functions
    to the Riemann sphere (which is another reason why
    zeroes and poles behave dually to one another).
    We will say more about meromorphic functions momentarily.

    \ii Finally, in the case where $ z_0$
    is neither removable nor a pole, we say that $z_0$ is
    an \vocab{essential} singularity ---
    these have incredibly interesting behavior.
\end{enumerate}

Here is one example of the crazy behavior of essential singularities.
(In my head, the situation is just that
cases (1) and (2) are so rigid and well-behaved
that deviating at all from those scenarios
causes things to explode. Maybe the right adjective is ``brittle''.)
\begin{prop}
    [Casorati-Weierstrass]
    \label{prop:essential-singularity-density}
    If $z_0$ is an essential singularity of $f$,
    then the image of any neighborhood of $z_0$ under $f$
    is dense in $\CC$.
\end{prop}

\begin{proof}
    \todo{todo; idea is just that if it isn't dense
    then we get a rem sing or pole}
\end{proof}

For example, $f(z) = e^{1/z} = \sum_{n=0}^\infty \frac 1{n!z^n}$
has an essential singularity at zero.
\\

Now, we will talk about meromorphic functions.
The most important definition relating to meromorphic functions
is the idea of a residue.

\begin{definition}[Residue]
    \label{def:residue}
    If $f$ is holomorphic in a punctured neighborhood of $a\in\CC$
    (so not necessarily holomorphic at $a$ itself), then
    the \vocab{residue} of $f$ at $a$ is
    \begin{equation*}
	\res_a(f) = \frac 1{2\pi i}\int_{\abs{z-a} = \eps} f(z)\dd z
    \end{equation*}
    where $\eps$ is small.
\end{definition}

Of course, residues are only interesting at singularities,
since if $f$ is holomorphic at $a$ then
the residue is zero by Cauchy's integral theorem
(Theorem \ref{thm:cauchy-integral-theorem}).

\begin{prop}
    \label{prop:residue-formula}
    Suppose $a$ is a pole of $f$ with order $n$. Then
    \begin{equation*}
	\res_a(f) = \frac 1{(n-1)!}\frac{\dd^{n-1}}{\dd z^{n-1}}
	((z-a)^nf(z))\Big|_{z=a}.
    \end{equation*}
\end{prop}

\begin{proof}
    By definition of pole, $(z-a)^nf(z)$ extends to
    a nonzero holomorphic function in a neighborhood of $a$,
    so there exist complex numbers $c_k$ such that
    \begin{equation*}
	(z-a)^nf(z) = c_0 + c_1(z-a) + c_2(z-a)^2 + \cdots,
    \end{equation*}
    with $c_0\ne 0$, in a neighborhood of $a$.
    This means that
    \begin{equation*}
	f(z) = \frac{c_0}{(z-a)^n} + \frac{c_1}{(z-a)^{n-1}}
	+ \cdots + \frac{c_{n-1}}{z-a} + c_n + c_{n+1}(z-a) + \cdots
    \end{equation*}
    in a punctured neighborhood of $a$.
    Integrating over a sufficiently small circle centered at $a$
    kills every addend except $\frac{c_{n-1}}{z-a}$
    (where we can integrate termwise because
    power series converge locally uniformly), so that
    \begin{equation*}
	\res_a(f) = \frac 1{2\pi i}\int_{\partial B_\eps(a)} f(z)\dd z
	= \frac 1{2\pi i}\int_{\partial B_\eps(a)}\frac{c_{n-1}}{z-a}\dd z
	= c_{n-1}.
    \end{equation*}
    On the other hand we can compute $c_{n-1}$ in the standard way
    by differentiating the power series for $(z-a)^nf(z)$
    above $n-1$ times, dividing by $(n-1)!$, and evaluating at $a$.
\end{proof}

In the case where $n = 1$, the formula above involves
the zeroth derivative of $(z-a)f(z)$, which is just
the value of (the extension to a neighborhood of $a$ of)
$(z-a)f(z)$ at $a$, i.e.
\begin{equation*}
    \res_a(f) = \lim_{z\to a}(z-a)f(z).
\end{equation*}
Such a pole is called a \vocab{simple pole} of $f$.

The reason that we care to define residues is that
the integral of a meromorphic function is dependent
only on the residues:
\begin{theorem}[Residue theorem]
    \label{thm:residue-theorem}
    Let $f$ be a meromorphic function defined on good domain $\Omega$.
    Then, $f$ has finitely many poles $a_1$, \dots, $a_N$, and
    \begin{equation*}
	\frac 1{2\pi i}\int_{\partial\Omega}f(z)\dd z
	= \sum_{k=1}^N\res_{a_k}(f).
    \end{equation*}
\end{theorem}

\begin{proof}
    Since the poles of $f$ are the zeroes of $1/f$
    (or rather, we can extend $1/f$ to a meromorphic function
    whose zeroes are the poles of $f$), the set of poles is closed,
    so it is compact and discrete, hence finite.

    In particular, there exists $\eps>0$ such that
    the $\eps$-balls around each $a_k$ are pairwise disjoint,
    meaning $f$ is holomorphic on
    $\tilde\Omega = \Omega - \bigcup_{k=1}^NB_\eps(a_k)$, so that
    \begin{equation*}
	0 = \int_{\partial\tilde\Omega}
	f(z)\dd z = \int_{\partial\Omega}f(z)\dd z
	- \sum_{k=1}^N \int_{\partial B_\eps(a_k)} f(z)\dd z
	= \int_{\partial\Omega}f(z)\dd z
	- \sum_{k=1}^N \res_{a_k}(f).\qedhere
    \end{equation*}
\end{proof}

\todo{computation of real integrals with residue theorem}

One way to think about the argument principle
(Proposition \ref{prop:argument-principle})
is that holomorphic functions look like polynomials of degree $n$
in a neighborhood of a zero of order $n$,
which we can glean from the power series,
as higher-order terms are small in comparison to the lowest one.
This explains why the winding number of $f(\partial B_\eps(z_0))$
around $0$ is $n$ if $\eps$ is small enough
and $z_0$ is an isolated zero of order $n$.
(Green's theorem lets us patch this local interpretation
into a more global integral, in the same way that
the residues of a meromorphic function determine its integral.)

In analogy, a meromorphic function $f$ looks like
$(z-a_0)^{-n}$ around a pole $a_0$ of order $n$,
which implies that the winding number of $f(B_\eps(z_0))$
for small $\eps$ should be $-n$
(since for example $1/z$ reverses the orientation of the circle).
Thus, the poles are ``negative zeroes'' for the purposes
of counting facts like the argument principle:
\begin{prop}[Argument principle for poles]
    \label{prop:argument-principle-poles}
    Let $f$ be meromorphic on a good domain $\Omega$
    with no zeroes or poles on $\partial\Omega$. Then
    \begin{equation*}
	\frac 1{2\pi i}\int_{\partial\Omega}\frac{f'(z)}{f(z)}\dd z
	= N(f,\Omega) - P(f,\Omega)
    \end{equation*}
    where $N(f,\Omega)$ and $P(f,\Omega)$ are respectively
    the number of zeroes and poles of $f$ in $\Omega$.
\end{prop}

Note that the integral on the left gets negated
when replacing $f$ with $1/f$, which makes sense
because taking the reciprocal swaps zeroes and poles.

\begin{proof}
    [Proof of Proposition \ref{prop:argument-principle-poles}]
    Note that $n = N(f,\Omega)$ and $p = P(f,\Omega)$ are both finite.
    We do a similar argument to the proof of
    Proposition \ref{prop:argument-principle}:
    if $a_1$, \dots, $a_n$ are the zeroes of $f$
    and $b_1$, \dots, $b_p$ are the poles, both lists
    with repeats according to multiplicity, then
    \begin{equation*}
	f(z) = \frac{(z-a_1)\cdots(z-a_n)}{(z-b_1)\cdots(z-b_p)}g(z)
    \end{equation*}
    where $g$ extends to a nowhere zero holomorphic function on $\Omega$.
    Then, by the product rule or logarithmic differentiation,
    \begin{equation*}
	\frac{f'(z)}{f(z)} = \sum_{k=1}^n\frac 1{z-a_k}
	- \sum_{k=1}^n \frac 1{z-b_k} + \frac{g'(z)}{g(z)},
    \end{equation*}
    and then taking the integral gives
    \begin{equation*}
	\frac 1{2\pi i}\int_{\partial\Omega}\frac{f'(z)}{f(z)}\dd z
	= \sum_{k=1}^n 1 - \sum_{k=1}^p 1 + 0 = n - p
    \end{equation*}
    by the Cauchy integral formula for the $\frac 1{z-w}$ summands
    and the Cauchy integral theorem for $\frac{g'(z)}{g(z)}$.
\end{proof}

Recall Exercise \ref{exer:rouche-zeroes} tells us how
if two holomorphic functions $f$ and $g$ send $\partial\Omega$
sufficiently close to one another (in the sense of the given inequality),
then $f$ and $g$ have the same number of zeroes in $\Omega$.
Since that version of Rouch\'e's theorem relied on
only the argument principle, there is a version for meromorphic functions:
\begin{theorem}[Rouch\'e's theorem]
    \label{thm:rouches-theorem}
    Let $f$ and $g$ be meromorphic functions on $\Omega$.
    If $\abs{f-g}<\abs f$ on $\partial\Omega$, then
    $N(f,\Omega) - P(f,\Omega) = N(g,\Omega) - P(g,\Omega)$.
\end{theorem}

The homotopy idea in Exercise \ref{exer:rouche-zeroes}
can also be used here, but here is another proof.

\begin{proof}
    [Proof of Rouch\'e's theorem]
    The inequality says that $\smabs{\frac gf - 1} < 1$
    on $\partial\Omega$, so $F = \frac gf$ maps $\partial\Omega$
    into the open disk $B_1(1)$. Let
    \begin{equation*}
	n(a) =\frac 1{2\pi i}\int_{\partial\Omega}\frac{F'(z)}{F(z)-a}\dd z
    \end{equation*}
    be the number of times $F$ attains the value
    $a\in\CC - F(\partial\Omega)$ in $\Omega$,
    which is constant on each component of $\CC - F(\partial\Omega)$.
    Obviously $n(a)$ is zero on the unbounded component
    (since e.g. $\smabs{\frac gf}$ is bounded on $\partial\Omega$),
    and $0$ is contained in the unboundded component,
    so number of zeroes of $F$ is
    \begin{align*}
	0
	&= n(0) = \int_{\partial\Omega}\frac{F'(z)}{F(z)}\dd z
	= N\left(\frac gf,\Omega\right) - P\left(\frac gf,\Omega\right)\\
	&= (N(g,\Omega) + P(f,\Omega))
	- (P(g,\Omega) + N(f,\Omega)),
    \end{align*}
    where we recall that taking the reciprocal swaps poles and zeroes.
\end{proof}


\todo{laurent series}

\subsection{Schwarz lemma}
The maximum principle (and the open mapping theorem)
are very strong structural statements,
and with some cleverness we can extract a lot of juice out of them.
\begin{theorem}[Schwarz lemma]
    \label{thm:schwarz-lemma}
    Let $\Delta$ be the open unit disk
    and $f\in\mathcal O(\Delta)$ such that
    \begin{enumerate}[(1)]
	\ii $\abs{f} < 1$ on $\Delta$; and
	\ii $f(0) = 0$.
    \end{enumerate}
    Then $\abs{f(z)} < \abs{z}$ on all of $\Delta$
    and $\abs{f'(0)} \le 1$. Furthermore, if $\abs{f(z)} = \abs z$
    for any nonzero $z\in\Delta-\{0\}$ or $\abs{f'(0)} = 1$,
    then $f(z) = cz$ for some $c\in S^1$.
\end{theorem}

\begin{proof}
    Let $g(z) = \frac{f(z)}{z}$.
    Since $f(0) = 0$, $f$ has a power series
    expansion $f(z) = c_1z + c_2z^2 + \cdots$
    that converges on all of $\Delta$
    (examine the proof of holomorphicity implying analyticity,
    Theorem \ref{thm:holomorphic-implies-analytic}).
    Define $g(z) = c_1 + c_2z + \cdots$
    to be the extension of $\frac{f(z)}z$ to $\Delta$
    by setting $g(0) = f'(0)$; this function is holomorphic
    because we have defined it via a power series.

    Given $z\in\Delta$, we have by the maximum principle
    (Theorem \ref{thm:max-mod-principle}) that
    \begin{equation*}
	\abs{g(z)} \le \sup_{w\in B_r(0)} \frac 1{\abs w}
	\le \frac 1r
    \end{equation*}
    for any $r$ satisfying $\abs z \le r < 1$,
    so taking $ r$ arbitrarily close to $1$ tells us that
    $\abs g \le 1$, which is equivalent to the inequalities
    that we want to prove. The two equality conditions
    are both equivalent to saying that $\abs{g(z)} = 1$
    for some $z\in\Delta$; this would imply that
    $z$ is a local maximum for $g$ in the interior of
    $B_r(0)$ for any $r$ with $\abs z < r < 1$,
    which implies that $g$ must be constant on $B_r(0)$.
    If this is the case then $g(0) = c_1$ has magnitude $1$
    and $f(z) = c_1z$ for all $z\in\Delta$, as desired.
\end{proof}

\begin{corollary}
    \label{coro:automorphism-disk}
    The automorphisms of the unit disk
    (i.e. holomorphic bijections $\Delta\to\Delta$ with holomorphic inverse)
    are given by the fractional linear transformations
    $\varphi(z) = c\frac{a-z}{1-\ol az}$
    for $\abs a < 1$, $\abs c = 1$.
\end{corollary}

\begin{proof}
    Let $a\in\Delta$ be the point with $f(a) = 0$,
    $\varphi(z) = \frac{a-z}{1-\ol az}$,
    and $g(z) = f(\varphi(z))$. Then, $g$ and $g\inv$
    are automorphisms of $\Delta$ and satisfy the hypotheses
    of the Schwarz lemma, so $\abs{g'(0)}\le 1$
    and $\abs{(g\inv)'(0)}\le 1$.
    The chain rule tells us that $g'(0)(g\inv)'(0) = 1$,
    so the inequalities above are equalities,
    meaning $g(z) = cz$ for some $\abs c = 1$.
\end{proof}

\todo{weierstrass preparation theorem (not required content)}

\pagebreak
\section{Harmonic functions}
\begin{definition}[Harmonic function]
    \label{def:harmonic-function}
    Let $U\subseteq\CC$ be open and connected.
    A $C^2$ function $u\colon U\to\RR$ is \vocab{harmonic} if
    \begin{equation*}
	\Delta u\coloneq u_{xx} + u_{yy}
	= 4\frac{\partial^2}{\partial \ol z\partial z} u
	= 0.
    \end{equation*}
\end{definition}
\todo{maybe write definition wrt $\RR^n$ instead of $\CC$}

\begin{example}
    Real parts and imaginary parts of holomorphic functions are harmonic
    by the Cauchy-Riemann equations, since if $f = u + iv\in\mathcal O(U)$,
    \begin{equation*}
	u_{xx} + u_{yy} = (u_x)_x + (u_y)_y
	= v_{yx} - v_{xy} = 0
    \end{equation*}
    and similarly for $v$.
    Thus, for example, $u(x,y) = e^{x}\cos(y)$,
    being the real part of $e^{x+iy}$, is harmonic.
    We will shortly see that all harmonic functions arise in this way.
\end{example}

Writing $\Delta u$ in terms of $\frac\partial{\partial z}$
and $\frac{\partial}{\partial\ol z}$ tells us that
$u\in C^2(U)$ is harmonic if and only if
\begin{equation*}
    \frac{\partial u}{\partial z}
    = \frac{\partial u}{\partial x} - i\frac{\partial u}{\partial y}
\end{equation*}
is harmonic (since applying $\frac\partial{\partial\ol z}$ is zero).

Let $U\subseteq\CC$ be a simply connected open set
and suppose that $u\in C^2(U)$ is harmonic. For a fixed $z_0\in U$,
\begin{equation*}
    f(z) = \int_{z_0}^z\frac{\partial u}{\partial \zeta}\dd\zeta
\end{equation*}
is a well-defined holomorphic function on $U$
(where the integral is along some arbitrary path from $z_0$ to $z$ in $U$).
Then,
\begin{equation*}
    \ol{f(z)} = \int_{z_0}^z \frac{\partial u}{\partial\ol\zeta}\dd\ol\zeta.
\end{equation*}
Recalling that $\dd u = \frac{\partial u}{\partial\zeta}\dd\zeta +
\frac{\partial u}{\partial\ol\zeta}\dd\ol\zeta$,
\begin{equation*}
    2\opname{Re}f(z) = f(z) + \ol{f(z)}
    = \int_{z_0}^z \dd u = u(z) - u(z_0).
\end{equation*}
Thus $u$ is (up to a constant) equal to the real part of $f$
(which justifies the remark at the end of the example above).
Let $u^*$ be the imaginary part of $f$, which is also a harmonic function,
called the \vocab{conjugate harmonic function} of $u$.
Note that if $g$ is another holomorphic function on $U$
whose real part is also $\frac u2$ up to a constant, then
$f-g$ is a holomorphic function whose real part is constant;
the open mapping theorem then implies that $f-g$ is constant,
so $u^*$ is the unique harmonic function up to a constant
such that $u + iu^*$ is holomorphic.

Harmonic functions inherit some nice properties of harmonic functions.
\begin{prop}[Mean value property]
    \label{prop:mean-value-thm}
    Let $u$ be harmonic on open $U\subseteq\CC$ with $z_0\in U$ a point.
    Then, if $U$ contains the disk $\ol{\Delta_R(z_0)}$,
    \begin{equation*}
	u(z_0) = \frac 1{2\pi}\int_0^{2\pi} u(z_0 + re^{i\theta})\dd\theta
    \end{equation*}
    for all $0 < r \le R$.
\end{prop}

\begin{proof}
    This follows immediately from using the Cauchy integral formula
    on $u + iu^*$ and comparing the real and imaginary parts of both sides.
\end{proof}

\pagebreak
\begin{prop}[Max-min principle]
    \label{prop:max-min}
    If $u$ is harmonic on open, connected $U\subseteq\CC$ and attains
    a local minimum or maximum, it is constant.
\end{prop}

\begin{proof}
    Let $z_0\in U$ be a local maximum for $U$.
    For $0 < r \le R$, where $R$ is small enough that
    the disk centered at $z_0$ with radius $R$ is contained in $U$,
    \begin{equation*}
	\sup_{\abs{z-z_0} = r} u(z) \le u(z_0)
	= \int_0^{2\pi} u(z_0 + re^{i\theta})
	\le \sup_{\abs{z-z_0} = r} u(z),
    \end{equation*}
    so $u(z_0)$ is equal to $\sup_{\abs{z-z_0}=r}u(z)$.
    Since $u$ is continuous, this tells us that
    $u$ is constant on $\{z\in U : \abs{z-z_0} = r\}$
    for all $r < R$, hence is locally constant.
    Since $U$ is connected, this implies $u$ is constant.
    A similar proof works for if $z_0$ is a local minimum.
\end{proof}

By taking an appropriate M\"obius transformation,
we can use the mean value property
to obtain the value of a harmonic function
on any point on the interior of a disk
given only the boundary values.
\begin{theorem}[Poisson formulas]
    \label{thm:poisson-formulas}
    Let $u$ be a continuous function on the closed disk
    $\ol\Delta$ and harmonic on $\Delta$. Then for any $a\in\Delta$,
    \begin{equation*}
	u(a) = \frac 1{2\pi}\int_0^{2\pi}
	\frac{1 - \abs a^2}{\abs{e^{i\psi} - a}^2} u(e^{i\psi})\dd\psi
	= \frac 1{2\pi}\int_0^{2\pi}
	\opname{Re}\left(\frac{e^{i\psi + a}}{e^{i\psi-a}}\right)
	u(e^{i\psi})\dd\psi.
    \end{equation*}
\end{theorem}
\todo{proof (annoying calculation)}

\todo{schwarz theorem}

\todo{reflection principle}

\pagebreak
\section{Todo items}
\listoftodos












\end{document}
